\documentclass[runningheads]{llncs}
%
\usepackage[T1]{fontenc}
% T1 fonts will be used to generate the final print and online PDFs,
% so please use T1 fonts in your manuscript whenever possible.
% Other font encondings may result in incorrect characters.
%
\usepackage{graphicx}
% Used for displaying a sample figure. If possible, figure files should
% be included in EPS format.
%
% If you use the hyperref package, please uncomment the following two lines
% to display URLs in blue roman font according to Springer's eBook style:
%\usepackage{color}
%\renewcommand\UrlFont{\color{blue}\rmfamily}
%\urlstyle{rm}
\usepackage{comment}
\usepackage{diagbox}
\usepackage{url}
%\usepackage{breakurl}
\usepackage[breaklinks]{hyperref}

\usepackage[utf8]{inputenc}
\usepackage[T1]{fontenc}
\usepackage{soul}
\usepackage{graphicx}
\usepackage{amsmath}
\usepackage{booktabs}
\usepackage{pifont}
\usepackage{enumitem}
\usepackage{tikz}
\usepackage{cite}
\usepackage{amsmath,amssymb,amsfonts}
\usepackage{algorithmic}
\usepackage{graphicx}
\usepackage{textcomp}
\usepackage{xcolor}
\usepackage{amsmath}
\def\BibTeX{{\rm B\kern-.05em{\sc i\kern-.025em b}\kern-.08em
    T\kern-.1667em\lower.7ex\hbox{E}\kern-.125emX}}
    
\def\checkmark{\tikz\fill[scale=0.4](0,.35) -- (.25,0) -- (1,.7) -- (.25,.15) -- cycle;} 

\usetikzlibrary{trees,arrows}

\newlist{todolist}{itemize}{2}
\setlist[todolist]{label=$\square$}
\newcommand{\cmark}{\ding{51}}%
\newcommand{\xmark}{\ding{55}}%
\newcommand{\done}{\rlap{$\square$}{\raisebox{2pt}{\large\hspace{1pt}\cmark}}%
\hspace{-2.5pt}}
\newcommand{\wontfix}{\rlap{$\square$}{\large\hspace{1pt}\xmark}}

\def\UrlBreaks{\do\/\do-}

\renewcommand\thesection{\arabic{section}} % arabic numerals for the sections


% correct bad hyphenation here
\hyphenation{op-tical net-works semi-conduc-tor IEEEconf}

\begin{document}
%% Example cover page for a two-page IEEE conference document
%\onecloumn
% Cover page

\begin{center}
  \textbf{recovery testbed}
\end{center}

\begin{itemize}
  \item DEADLINE: Abstract April 11, Paper April 18, 2025
  \item Publication Venue: QEST+FORMATS 2025
  \item Url for Publication Venue: \url{https://www.qest-formats.org/}
  \item Submission Category or Track:
    \begin{itemize}
        \item Models and metrics for the correctness, performance, reliability, safety, and security of systems (stochastic, probabilistic, and non-deterministic models including Markov chains, automata, Petri nets, process algebra, max-plus algebra, and their variations);
    \end{itemize}
  \item Page limit:  16 pages; currently at 16
  \item Research Contribution: 
   The proposed model enables comparisons of alternative
repair schemes by considering the availability of repair crews,
predicting failure propagation, and estimating the level of
service capacity over time. The model can also guide staffing
decisions by identifying the minimum number of repair crews
required to maintain a particular service level.
\end{itemize}

\vspace{2ex}

%SSS- This helps us keep track of the number of drafts
\textbf{Version 1.1}
Last Updated: 4/1/2025 \\
\textbf{Major Changes Since Previous Draft}
\begin{itemize}
    \item Case study
    \item Conclusion
\end{itemize}

\vspace{2ex}

\textbf{Outline}
\begin{enumerate}
%SSS - Include the number of pages in your current draft in front of each section (not subsection) title. Also include the status (complete, pending, not written).
  %SSS - The following sections should be in every paper. 
  \item  Abstract
  \item Introduction 
  \begin{enumerate}
          \item Blackout in Kenya
          \item Motivation for this paper
          \item Intro to this model
        \end{enumerate}
  \item Related work
  \begin{enumerate}
          \item Modeling of cascading failure
          \item Modeling of the Recovery process
        \end{enumerate}
  \item Proposed recovery model
  \begin{enumerate}
          \item Modeling of Cascading failure
          \begin{enumerate}
              \item States
              \item Actions
              \item Rewards
              \item Transition probabilities
          \end{enumerate}
          \item Repair schemes
          \begin{enumerate}
              \item Randomized priorities
              \item pre-made repair priorities
              \item Look-ahead dynamic priorities
              \item cascading effect dynamic priorities
          \end{enumerate}
        \end{enumerate}
  \item case studies
  \begin{enumerate}
          \item IEEE-14
          \item Simulation Environment
          \item IEEE-14 Failure Cases
          \begin{enumerate}
            \item experiment 1
            \item experiment 2
            \item Experiment 3
            \end{enumerate}
        \item Discussion of results
        \end{enumerate}
  \item Conclusion and Future Work 
    \begin{enumerate}
          \item The contribution to the body of knowledge 
          \item social impact 
          \item future directions, enable advances in the repair scheme area
        \end{enumerate}
  \item Acknowledgements 
  %SSS - Include this text in your acknowledgments: The authors gratefully acknowledge funding from NSF DUE-1742523, DEd GAANN awards P200A180095 and P200A210121, and the Missouri S\&T Intelligent Systems Center.
  \item References
\end{enumerate}

\vspace{2ex}

\textbf{TO DO (Advisors):}
\begin{itemize}
    \item review the paper
\end{itemize}

\textbf{TO DO (Sanfan):}
\begin{itemize}
    \item review the paper entirely
    \item submit for publication 
\end{itemize}

\textbf{Author guidelines (submission requirements)}: \\ 
\url{https://www.qest-formats.org/cfp.html}\\
Regular papers must not exceed 16 pages, excluding references.
We are considering additional presentation-only submission types; details will be provided in the later CfPs.

All papers must be submitted in Springer’s LNCS format and will undergo a rigorous single-blind review process. All submitted papers must be unpublished and not be submitted for publication elsewhere. Papers should be submitted electronically using EasyChair at this link.

Springer encourages authors to include their ORCIDs in their papers. Authors should consult Springer’s authors’ guidelines and use Springer’s LaTeX templates for the preparation of their papers. Submitted papers not complying with the above guidelines may be rejected without undergoing review.
We are considering awards for best papers and best artifacts.

Publications\\
All accepted papers need to be presented and discussed at the conference by one of the authors. The QEST+FORMATS 2025 proceedings will be published in the Springer LNCS series indexed by ISI Web of Science, Scopus, ACM Digital Library, dblp, Google Scholar. All submitted papers will be evaluated by at least three reviewers on the basis of their originality, technical quality, scientific or practical contribution to the state of the art, methodology, clarity, and adequacy of references.

Special Issue\\
We are working on a special issue in a Q1 journal. A selection of the best papers will be invited to submit an extended version of their work to special issues in internationally recognized journals.

Artifact Evaluation\\
Reproducibility of experimental results is crucial to foster an atmosphere of trustworthy, open, and reusable research. To improve and reward reproducibility, QEST+FORMATS 2025 will include a dedicated Artifact Evaluation (AE). Submission of an artifact is mandatory for tool papers (both regular and short), and optional but encouraged for research and case study papers where it can support the results presented in the paper. More details can be found in the corresponding page
%End Cover Sheet



%SSS - Need a better title. Something like: "Win the race: Decision support for recovering complex systems as failures propagate"
%SSS - Maybe something with an acronym that amounts to "RACE"
\title{A recovery process model for complex networked system under cascading failure}
%
%\titlerunning{Abbreviated paper title}
% If the paper title is too long for the running head, you can set
% an abbreviated paper title here
%
\author{Sanfan Liu\orcidID{0000-0002-5089-3500} \and
Sahra Sedigh Sarvestani\orcidID{0000-0002-2917-5643 } \and
Ali R. Hurson\orcidID{0000-0002-9510-0762}}
%
\authorrunning{S. Liu et al.}
% First names are abbreviated in the running head.
% If there are more than two authors, 'et al.' is used.
%
\titlerunning{A with time-inhomogeneous transitions} %change page header title (short title)

\institute{Missouri University of Science and Technology \\
\email{\{slbdb,sedighs,hurson\}@mst.edu}}
%
\maketitle              % typeset the header of the contribution
%
\begin{abstract}
%SSS - add at least a bulleted outline  
\begin{itemize}
    \item Cascading failure has been a great hazard for complex networked systems  (very brief)
    \item Related work are either dependency analysis with expansive modifications without guarantee to work, or recovery model with too much limitation or unrealistic assumptions (brief)
    \item A TI-MDP is proposed in this paper to analyze the recovery process of a complex networked system (be as specific as possible)
    \item IEEE-30 and 14 are the case study done to compare some common recovery schemes found in related works
    \item The TIMDP was used to find the most suitable schemes and enables future research in recovery schemes (b/w this bullet, and the last one, you create lack of clarity about whether any new schemes are being proposed; if a fundamental use of this work is comparison based on specific criteria of schemes, that should be very clear 
\end{itemize}
\end{abstract}

%SSS - an attractive figure can hook in your reader and tell them that you know exactly what you are doing
\section{Introduction}
Cascading failures represent a fundamental challenge in the study of interdependent and multilayered infrastructures. The malfunction of a single component can trigger a chain of disruptions that propagate swiftly across interconnected subsystems. Such dynamics are particularly critical in domains such as water distribution, urban transportation, and smart grids, where the coupling of physical distribution layers with cyber-communication networks heightens systemic fragility.

The socio-technical consequences of cascading failures are severe. Recovery trajectories are often prolonged, economic costs escalate rapidly, and extended service interruptions threaten societal stability. Addressing these risks requires analytical frameworks that capture the nonlinear dynamics of failure propagation while also accounting for the contingent role of human and organizational decision-making. This paper engages with these challenges by examining the mechanisms of cascading failure, critically analyzing past responses, and presenting strategies to enhance resilience in complex networked systems.

A productive way to conceptualize cascading failure is as a race against propagation. Once a perturbation occurs, the pace of escalation accelerates as more components fail. Effective intervention therefore depends not only on rapid response but also on identifying and prioritizing structurally critical nodes. Poorly conceived or random repair strategies—as illustrated by the Kenyan blackout—waste valuable time while allowing the cascade to intensify. In contrast, interventions guided by rigorous modeling can reduce disruption and accelerate recovery.

%SSS - mention the year right here so people know it is a very recent example
The nationwide blackout in Kenya underscores the detrimental consequences of ineffective recovery decisions. A miscommunication at the Lake Turkana wind facility triggered the sudden loss of 270 MW, roughly 15\% of national demand. Attempts to stabilize the grid through imports from Uganda proved ineffective, delaying restoration and extending the outage to nearly 50 million people ~\cite{apkenya_aug23, apkenya_dec23, mwangi_kenya_2023}. In contrast, the 2025 Iberian Peninsula blackout highlights the value of coordinated, knowledge-driven response. Operators activated cross-border interconnections with France and Morocco, implemented controlled load shedding, and preserved critical infrastructure, thereby limiting systemic damage and expediting recovery ~\cite{cetinkaya_technical_2025}. These contrasting cases illustrate how decision-making can either exacerbate or mitigate cascading failure.

%SSS - The emphasis on interdependency should serve as a lead-in to explaining that your work is one of few studies that explicitly considers it 
Recent research has emphasized the central role of interdependencies in driving fault propagation. Studies consistently demonstrate that tightly coupled architectures amplify systemic fragility, while strategic modifications—such as redundancy, capacity enhancement, or load redistribution—can improve survivability by slowing or halting cascades~\cite{zhang_effects_2017}. Yet, such measures are not always feasible due to environmental constraints and the inherent complexity of real networks. Even advanced smart grids, as shown in the Iberian case, remain vulnerable to cascading collapse, underscoring that recovery planning often constitutes the final line of defense against catastrophic failure.

Parallel scholarship has explored recovery modeling, but many approaches rely on restrictive assumptions about repair times and conditions. Some treat recovery durations as fixed~\cite{sarkale_solving_2018,biswas_repair_2025}, while others assume exponential distributions~\cite{yan_dynamic_2021}. Fixed durations may be adequate when failures are measured in days, yet they obscure the finer temporal scales at which cascading failures unfold, often within hours or minutes. Exponential assumptions are somewhat more realistic at these shorter timescales, but they fail to capture the heterogeneity consistently observed in practice. Empirical studies frequently report Lognormal~\cite{l_c_cadwallader_review_2012} or Weibull~\cite{ocansey_temperature-driven_2025} distributions, underscoring the need for models that flexibly represent diverse recovery dynamics. 

Furthermore, many recovery modeling studies presuppose that repairs occur only when the system is effectively immune to additional failures—either because it has been shut down or because isolation mechanisms prevent further propagation~\cite{sarkale_solving_2018, yan_dynamic_2021,wang_q-learning-based_2025}. This assumption is unrealistic for critical infrastructures, where restoration efforts typically proceed under ongoing risk. The blackout in Kenya, among other rolling blackouts, serves as a cautionary example of the dangers of neglecting this reality.

%SSS - Model for what? This sentence is critical. Make it boldface. 
Building on these insights, the objective of this study is to present a mathematical model with broad applicability. The model provides predictive insights on recovery strategies for complex networked systems, deepening understanding of the race between cascading failure and restoration. It supports improved network design, optimized task scheduling, and efficient allocation of recovery resources. For decision makers, it enables proactive planning by identifying robust recovery strategies in advance of disruptions, thereby establishing effective default protocols. During real-time cascading events, the model supports adaptive policy adjustments, allowing operators to recalibrate strategies dynamically and mitigate disruption before the race against failure propagation is irretrievably lost.

Earlier approaches often assumed that repair times were either known in advance or followed an exponential distribution, but such assumptions rarely hold across the diverse contexts of real-world infrastructure. To address this, the present model incorporates time-inhomogeneous transitions, which accommodate a wider range of empirical repair time distributions and reflect the temporal heterogeneity of restoration processes.

Within the Markov framework, future transitions depend only on the current state and are independent of prior states. Yet the practical effort invested in repairing failed components cannot be erased each time the system state evolves, nor can repair be treated as instantaneous. To reflect this reality, time-inhomogeneous transitions are employed to track the duration of repair activity for each component during an ongoing cascade. This approach enables the incorporation of distributions such as Lognormal or Weibull, thereby extending the model’s applicability to settings where exponential assumptions are inadequate.

Accordingly, this paper proposes a Markov decision process (MDP) model with time-inhomogeneous transitions that  reflect recovery in complex networked systems subject to cascading failure. The model further reflects dependencies between physical and cyber components as drivers of cascade dynamics. In earlier work \cite{marashi_identification_2021}, these interdependencies were identified and quantified using correlation metrics with a heuristic causation analysis method, which represent the conditional probability of failure propagation between interconnected components. When a component initially fails, it poses the greatest threat to system stability by generating error messages or faults to connected elements. As time progresses, the influence of a failed component gradually diminishes as failed components shut down, or the system adapts to absorb the faults. The fault propagation probability therefore begins at its peak upon initial failure and declines throughout the recovery process until the component is restored, at which point propagation ceases. Time-inhomogeneous transitions enable the model to capture this temporal evolution by linking the system state to the moment of failure and replicating the trajectory of fault propagation.


%%%%%%%%%%%%%%%%%%%%%%%%%%%%%%%%%%%%%%%%%%%%%%%%%%%%%%%%%%%%%%%%%%%%%
%\begin{comment}

\section{Q\&As: (Questions in comments) }
%Kenya's blackout due to cascading failure without a good recovery scheme
For example, a cascading failure happened in Kenya \cite{apkenya_dec23} that took 14 hours to recover and affected 50 million people \cite{apkenya_dec23}. The triggering event was an error message sent to the Lake Turkana wind power plant causing it to go offline \cite{apkenya_dec23}, losing 270 MW in power generation, which is 15\% of the local demand. This triggering event led to an imbalance of power and shut down other main generators, causing a large blackout area in Kenya, so they had to import power from a nearby country to begin the recovery process. However, after they set up the connection with Uganda's power grid, there was no power coming from Uganda \cite{apkenya_aug23}. At the same time, the failure continued to propagate throughout the country, resulting in a nationwide blackout. This example is one of many that highlight the importance of the careful selection of recovery actions and the serious consequences of cascading failures, which pose serious challenges to the stability, sustainability, and survivability of complex networked systems \cite{xiaohui_review_2016}. 

%SSS - Comparison to the recent blackout in Spain would be a good idea
%SSS - explicitly mention the Spain blackout
the 2025 Iberian Peninsula (Spain) blackout illustrates how coordinated, informed decision-making can limit damage. Operators leveraged cross-border support from France and Morocco, applied controlled load shedding, and prioritized critical infrastructure, effectively containing the disruption and speeding restoration~\cite{cetinkaya_technical_2025}.

Grid following-inverter-based resources (GFL-IBR) was the root cause of the cascading failure, as it will automatically match the renewable energy generation with the power flow in the grid. When a power error occurs in the grid, all the GFL-IBR will try to match with the error, caused the generators to shut down.
If there are more Grid forming inverter-based resources (GFM-IBR) installed in the system, the cascade could be eliminated. GFM-IBR can generate the “correct”(or say a “fixed”) quality of power for the gird. If the GFL can follow the GFM then the cascading failure could be reduced.
However, GFM-IBR are expansive and complex to make, to adjust the setting for the best quality of power. It’s a new technology, and we don’t know much about its consequences for massive installation.


To predict failure propagation paths and make mitigation designs, many mathematical models have been developed in recent years to elucidate the effect of failures and estimate the rate of recovery. 

%SSS - Introduce the concept of winning a race against failure propagation, and explain that is what we to enable here. 
When a random fault is detected in a complex networked system, it does not take long for the fault to spread to the next component via coupling. And as more and more components become failed, the rate of spread will increase faster and faster. The repair crews can start working on recovering the system while cascading failure is happening, but they have to make the right decision. If they don't know what they are doing, like the story in Kenya, they spend hours recovering an interconnection that does not help the system; in the meantime, the failure already spreads nationwide. 
Therefore, to win the race against failure propagation, we cannot just randomly repair failed components, we must target the critical ones that are essential to the system. 

Our model provided an environment for testing the recovery strategy in a complex networked system under cascading failure. Our model was designed with flexibility for various types of systems, whether they're power, water, gas, or data networks. 
They could have different interdependencies between components, but we can quantify them based on previous failure cases. These interdependencies will be reflected in failure propagation captured by our model. 

They could have different failure scenarios that could be 1, 2, or more failures happening to the system selected randomly or in a meaningful way to test the system's survivability facing different types of hazards from small, as a minor error, to large, like a natural disaster. Our model is applicable regardless of the distribution of repair time of components. 

%SSS - explain how the model is useful to others. Why should they care?
%Others include other researcher, power utilities, government decision/policy makers, etc. Anything we say here needs to be firmly backed up in the text and reiterated in the conclusions. 
%SSS - the statement immediately below describes the simulation environment, not the model. We simulate only smart grids, not other complex networks. 
Our model is valuable in various applications; it gives prevision and strategy for recovering a complex networked system. It enables deeper understanding of the race between recovery and cascading failure. It allows better design, task scheduling, and efficient use of recovery resources. 

For decision makers of the complex networked system, they can use the model to find out the most suitable way to recover the network before anything has happened, the default rules. And they adjust the policy during a cascading failure to find the best way to mitigate it before they completely lose the race.


%SSS - add a very brief explanation of what is missing from the body of knowledge that makes our work necessary. For example, "nothing exists that would help people decide how to win the race". Don't forget to check the paper recommended by the QEST reviewer - the one where failures were happening along with recovery. 

Studies in recent years are mainly focused on cascading failure analysis or recovering a complex networked system without any cascade. 
%SSS - Sanfan means that analysis of cascading failure is not enough .
%SSS - explicitly mention the paper that proposes tweaking dependencies and explain that it's costly and potentially infeasible as a solution
Analyzing the causes of cascading failure could help network managers to enhance or tweak the dependencies between components at an expensive cost, and due to the complexity of the network, cascading failures could still happen, like the grid in Spain. 
The recovery model without considering cascading failure is not useful either, as most of the complex networked systems are serving as the essential service of the society; they cannot be switched off for a post-hazard recovery or wait till the cascading failure finishes to begin recovery. The systems must be recovering while the cascading failure, and the repair crews must win the race before the failure becomes a serious problem that leads to a catastrophic failure with heavy penalties. 

Now, there are some studies that focus on recovering a complex networked system while failures occur. However, they assumed the repair time of each component is previously known, or assumed the repair time was an exponential distribution. The formal assumption is possible if the time scale is large, like if we are talking about days. People can do a lot of things in one day, and they can definitely fix whatever was broken, but the large time scale could miss lots of details in a cascading failure. It's also impossible to use on a cascading failure that recovery needs to be done in hours. The later assumption is better for cascading failure in the hourly scale, or in minutes. But it's not always true; the recovery process of a component is not always in exponential distributions. Uniform distribution and Weibull distribution are also commonly seen in the recovery time of a component, and the model should be flexible to accept those distributions. %SSS - add a reference for uniform/Weibull repair time

%SSS - refer to a specific paper for each of these statements (including the ones above)
Other than that, the recent work in recovery models has overlooked the importance of other layers of the system. For example, when recovering a smart grid, the research usually focuses on recovering the physical layer of the network, ``trying to get as much line powered as possible,'' but overlooks the importance of other layers of the system, such as the cyber layer. The cyber layer, or the control system layer, describes the cyber components of the system that adjust the flow and stabilize the quality of the service, and is an important element to consider when recovering a failing system.

%SSS - add a very brief reason for choosing a TI MDP to create what is missing from the body of knowledge
As mentioned above, some research assumes that the repair time is known, or in an exponential distribution. And we have discussed that this is not always possible or true, and we must consider other types of distribution as well. Time-inhomogeneous transitions are introduced to this model to solve the issue with different repair time distributions. 

As a Markov model, the future transition only relies on the current state and is independent of the previous states. However, the effort to repair spent on the failed components cannot be undone each time the model transitions to a new state. Also, the repair cannot be done in one time step like a miracle, so the transition for the repair completion process must be tracked by the time each crew spent on the component. Therefore, time-inhomogeneous transitions are utilized in this model to track the repair process of each component during the cascading failure. With TI, other distributions can be put into the repair time.

%SSS - need to mention that the model reflects dependencies among components, by using our previous work
This paper presents a Markov decision process (MDP) model with time-inhomogeneous transitions that reflect the recovery process of a complex networked system under cascading failure. 

The model considers dependencies between physical and cyber components as the root cause of cascading. The identification and quantification of the interdependency were present in our previous work \cite{marashi_identification_2021}, and then we created a dependency index to represent the conditional probability of failure propagation between components. When a component goes down, it could generate error messages or faults to other components that it connects to. When a component just goes down, it could bring the most impact to the system. As time continues, the impact will become less and less as the failed components would shut down or the system can bear the faults. Therefore, the fault propagation probability is going to start at the highest value when the component fails and decrease throughout the rest of the recovery process until the component gets restored and the fault propagation will no longer happen. The time-inhomogeneous transition can be used with the system state that tracks the failed moment of a component and replicates the fault propagation process. 

%\end{comment}
\section{Related work}

%SSS  - needs a bridge here to explain why we are considering these particular categories
To analyze the system during the cascading events, we must look at two particular topics. 
First, we must look at what is happening in the system, how the fault propagates, where the fault propagates to, and the time it takes for the fault to propagate; those questions go under the topic of cascading failure analysis. 
Then, we must know how the repair crews operate, how long does it take them to fix each component? how long does it take them to travel? How are the repair tasks assigned to the repair crews? Those questions go under the topic of recovery modeling. 


%SSS - try to identify the mostly closely related work and differentiate your work from it
Wang et al. proposed a cyber-physical power system model to test their Q-learning-based recovery method in \cite{wang_q-learning-based_2025}. Their model was designed to analyze the coupling between components within a CPPS and to prove their recovery method was superior to other common methods. 

%SSS - tactfully identify related work that has mistakes
The first outstanding mistake in their model is that recovery tasks can only begin after the cascading failure has finished and the system becomes stable. As mentioned previously, waiting for the cascading failure to end is unrealistic. Network managers must take action whenever the failure has been detected, and do their best to stop it from propagating, and the system may not be stable until it falls apart.

The second outstanding mistake is that each repair was used the exact same time (one time step) to repair. Even their model allows parallel repairs (so when two crews are assigned to repair any two different components, they will always be done at the same time).

The last one, in how they identify cascading failure. They assumed that if a component doesn't get any power from the generator, is considered "failed" and requires repair. They should only be "unpowered", as they are not broken, and they don't need repair to be back online. It's like you don't need to fix every power line after a blackout; that's a huge amount of work. 

\subsection{Related work in cascading failure}
Try to prevent and mitigate cascade failure, but it can be expensive. Our model focuses on the recovery side, which does not increase the cost.
There is so much uncertainty in complex networked systems, that with preventative and mitigative design, cascade failures can still happen. 

\subsection{Related work in recovery models}
Try to recover a system in a failed state, but some studies didn’t consider cascade failure (instead they assumed post-hazard, or isolated network system), and it could be an essential problem to the infrastructure. Some studies consider repair an “exponential distribution” which is not always true (made incorrect assumptions or contradictory assumptions on the repair time). The repair time could be in various distribution bases on different repair tasks. Our model has no limitation on the distribution of repairing a component.


\section{Methodology}

%SSS - all of your reasons for choosing a Markov model, MDP, TI, specific prob distr need to be crystal clear. Basically address the QEST comments with surgical precision. 

\subsection{Assumptions}
\begin{itemize}
    \item \textbf{Observability of Components:} All system components are considered fully observable to the decision maker, leveraging contemporary networked monitoring infrastructures. In contexts where direct observation is infeasible, system states can be inferred through sensor-generated outputs communicated to the control center, ensuring a high-fidelity representation of system status.
    \item \textbf{Repair Time Estimation:} The temporal dynamics of component restoration are modeled using empirical distributions derived from historical maintenance data, incident reports, or operational logs. These distributions accommodate variability due to failure severity, resource allocation, and environmental contingencies, providing a realistic and data-driven characterization of repair processes.
    \item \textbf{Failure Scope: } We choose to focus on the failures of transmission lines and smart control devices within the complex network, providing a focused and tractable modeling scope.
    %add reference to why they are more likely to fail
    \item \textbf{Post-Repair Reactivation: }Components are assumed to resume operation immediately after repair, with boot-up and testing activities included in the repair time, since these tasks require the  crew’s continued presence at the component.

\end{itemize}
\subsection{Markov chain model}
\subsubsection{System state}
Component state, repair crew states, ordered pairs for failed components \& failure time, ordered pairs for busy repair crews \& dispatch time.

\subsubsection{Time-inhomogeneous transitions}
\paragraph{Cascading failure}
    How to capture interdependency with TI
\paragraph{repair of a component}
    How to address general repair distribution with TI
\subsection{Decision support}
How are the repair tasks assigned to crews

\section{Case study}
\subsection{Repair scheme tested}
\paragraph{Randomized}
\paragraph{Pre computed}
\paragraph{Max-capacity}
\paragraph{Min-cascade}

\subsection{ Environment \& simulation software setup }
%SSS - mention that data sets, code, etc. are publicly available for reuse
Matlab, PSAT, python
\subsection{Results}
\paragraph{IEEE-14 for detailed inspection}
\paragraph{IEEE-30 for scalability of the model}
\section{Conclusion}
Towards a better decision support for complex networked systems with multi-layer dependency, which can cause cascading failures. A better chance to win the race against cascading failure.

We address some limitations and mistakes from the related works in the field and developed this TI-MDP model to compare the repair schemes.

\section{Future work}
%SSS - very clearly articulate the limitations, remember to mention cost of recovery
Identify limitations and ways to relax them
Other opportunities in the field of Deep learning, network performance analysis, etc. 

\bibliographystyle{IEEEtran}
\bibliography{export-data.bib}
%SSS - export-date is still blank

\end{document}
